
%\subsection{Overview of the Neural Shape Metrics Framework}
\subsection{Comparing neural populations as points in a metric space}


Consider a set of neural recordings collected across $K$ biological systems---for example, $K$ different animals, $K$ behavioral sessions from the same animal, or recordings from $K$ different brain regions (\textit{Fig.}~\ref{fig:schematic}a-c).
The activity of each neural system is recorded at $M$ measurement points, corresponding to different experimental conditions as well as different time points.
For example, neural responses could be measured in response to $C$ conditions (e.g. distinct sensory stimuli or trained motor actions) over $T$ time points per condition.
This would result in a total of $M = S T$ measurements.

Our high-level goal is to characterize the extent to which these $K$ neural systems are similar or different from each other across these measurement points.
To simplify the exposition, we will assume that each neural population is composed of $N$ neurons.
This assumption is often violated.
However, all the methods we outline can be trivially extended to the case of unequal neural population sizes (see \textit{Methods} and \cite{harvey2024duality} for mathematical proofs).


We will focus on comparing trial-averaged neural responses (see \textit{Discussion} for extensions that consider single-trial statistics).
In this setting, the data collected from $K$ neural systems are represented by a set of $N \times M$ matrices, $\mbX_1, \mbX_2, \dots, \mbX_K$.
Each row corresponds to a neuron's tuning curve the $M$ measurement points (task conditions and time points).
Each column corresponds to an $N$-dimensional population response to one of the experimental measurement.
\textit{Figure}~\ref{fig:schematic}d illustrates an example where the measurements span a 1-dimensional space, resulting in a 1-dimensional, smooth manifold structure within the neural firing rate space.

% As a simple first example, we will consider tasks that produce smooth, periodic tuning curves encoded by a task parameter $\theta \in [-\pi, \pi)$, which is discretized into $M$ bins.
% This is a generic form of tuning that pervades the neuroscience literature.
% Examples include representations of heading direction (e.g. \cite{chaudhuri,duszkiewicz}), orientation of a visual grating (e.g. \cite{hubel}), target locations in radial reaching tasks (e.g. \cite{georgopoulos}), and saccade targets in working memory tasks (e.g \cite{funahashi}).
% In all of these cases, the neural firing rates will trace a smooth ring manifold in $N$-dimensional firing rate space (\textit{Fig.}~\ref{fig:schematic}d).

Our first objective is to define a family of distance functions $d(\mbX_i, \mbX_j)$ that quantify the dissimilarity between two neural systems.
The principal challenge is that the neurons are randomly sampled and labeled in each system and so cannot be immediately compared.
This can be interpreted geometrically as two neural manifolds being randomly oriented with respect to each other in firing rate space (\textit{Fig.}~\ref{fig:schematic}e, top).
Although neural dimensions are not matched between systems, the $M$ measurements are shared (coloring in \textit{Fig.}~\ref{fig:schematic}d-e).
Thus, we can consider \textit{aligning} the two manifolds through some transformation and then computing some measure of distance (\textit{Fig.}~\ref{fig:schematic}e, bottom). 
% This is the motivating intuition for the shape metrics framework as well as several existing methods (e.g. linear predictivity scores and CCA).
% Intuitively, difference choices for the alignment transformation lead to different notions of distance between systems.

We will argue that it is useful to focus on distances which satisfy the following three axiomatic properties:
\begin{align}
\label{eq:metric-space-equivalence}
\text{Equivalence:} & \quad d\left(\mbX_i, \mbX_j\right)=0 \qquad \text{if and only if $\mbX_i$ and $\mbX_j$ are ``perfectly alignable''}\\
\label{eq:metric-space-symmetry}
\text{Symmetry:} & \quad d\left(\mbX_i, \mbX_j \right) =d\left(\mbX_j, \mbX_i \right)\\
\label{eq:metric-space-triangle}
\text{Triangle Inequality:} & \quad d\left(\mbX_i, \mbX_j \right) \leq d\left(\boldsymbol{X}_i, \mbX_k \right) + d\left(\mbX_k, \mbX_j \right)
\end{align}
where $\mbX_i$ and $\mbX_j$ are said to be ``perfectly alignable'' if the transform class specified by the practitioner is flexible enough to superimpose the two manifolds.

In mathematics, notions of distance that satisfy \cref{eq:metric-space-equivalence,eq:metric-space-symmetry,eq:metric-space-triangle} are said to define a \textit{metric space}.
Intuitively, metric spaces are useful because they enable us to spatially reason about data.
That is, we can envision each of our recordings from $K$ neural systems $\mbX_1, \dots, \mbX_K$ as points within the space of all possible neural recordings (\textit{Fig.}~\ref{fig:schematic}f).
That is, in this (potentially high-dimensional) metric space, an entire neural system is represented by a single point, and the distances between points indicate the geometric similarity or dissimilarity of the systems.   
These properties allow us to leverage methods such as clustering and nearest-neighbor regression (\textit{Fig.}~\ref{fig:schematic}g) and to do so with theoretical guarantees \parencite[see e.g.][]{Cover1967,Dasgupta2005}.

The \textit{Procrustes shape distance} \parencite{schonemann1966generalized} is a metric that satisfies \cref{eq:metric-space-equivalence,eq:metric-space-symmetry,eq:metric-space-triangle}, and most of our analyses will leverage this distance as a running example of the shape metrics framework.
If we assume that each neuron's tuning curve has been pre-processed to be mean centered (i.e. rows of $\mbX_i$ and $\mbX_j$ have mean zero), then the Procrustes distance can be written as:
\begin{equation}
\label{eq:procrustes-distance}
d_{\text{Proc}}(\mbX_i, \mbX_j) = \min_{\mbQ \in \cO(N)} \Vert \mbQ \mbX_i - \mbX_j \Vert_F
\end{equation}
Here, $\cO(N)$ is the set of $N \times N$ orthogonal matrices.
Intuitively, $\mbQ$ encodes the optimal rotation and reflection that minimizes the sum of squared residuals between the two neural manifolds.
The mean centering step ensures that the rotation and reflection is performed around the manifold's center of mass.

% Establishing a metric space enables us to use a broad set of off-the-shelf analysis tools.
% These methods take in a matrix of pairwise distances $\mbD_{ij} = d(\mbX_i, \mbX_j)$ and output unsupervised or supervised summaries of variability across the neural systems.
% For example, we may be interested in studying individual variability across $K$ animal subjects.
% To identify clusters of subjects with similar neural tuning, we may use hierarchical clustering methods, which work with theoretical guarantees in metric spaces~\parencite{Dasgupta2005}.
% To visualize these clusters, we can use multidimensional scaling (MDS) and other vector embedding techniques~\parencite{Borg2005-gk}.
% Further, if we measure a behavioral variable for each animal (e.g. performance on a task), we can use the distance relations between animals, based on neural data, to fit models that predict this behavior.
% For example, nearest nearest neighbor regression works with theoretical guarantees in metric spaces~\parencite{Cover1967}, and we will later demonstrate that it works well in practice.

% How can we define meaningful notions of distance between neural systems?
% This problem is nontrivial because simple measures, such as the Euclidean distance:
% \begin{equation}
% \Vert \mbX_i - \mbX_j \Vert_F = \left ( \sum_{n=1}^{N} \sum_{m=1}^{M} \left(\mbX_i[n, m] - \mbX_j[n, m] \right)^2 \right )^{1/2}
% \end{equation}
% are clearly meaningless whenever neurons are not identified and one-to-one matched across neural systems (\hl{Figure}).
% Intuitively, we need some way of \textit{aligning} the neural representations across two systems.
% We propose to use a constrained set of linear transformations, parameterized by a matrix $\mbQ$, to perform this alignment.
% Further, since it is common to pre-process neural data (e.g. by z-scoring each neuron) we propose to apply a user-defined function $\phi(\cdot)$ before fitting the alignment.
% Putting these pieces together, we can define:
% \begin{equation}
% \label{eq:generalized-euclidean-shape-metric}
% d(\mbX_i, \mbX_j) = \min_{\mbQ \in \cG} \Vert \mbQ \phi(\mbX_i) - \phi(\mbX_j) \Vert_F
% \end{equation}
% where $\cG$ denotes the set of allowable linear transformations.


% \begin{equation}
% \min_{\mbQ \in \reals^{N \times N}} \Vert \mbQ \mbX_i - \mbX_j \Vert_F \neq \min_{\mbQ \in \reals^{N \times N}} \Vert \mbX_i - \mbQ \mbX_j \Vert_F.
% \end{equation}

The Procrustes distance can be thought of as a constrained version of linear regression with an orthogonality constraint---i.e. a constraint such that the weight matrix must only encode a rotation and reflection of the firing rate space.
Na\"ively dropping this restriction is problematic for constructing metric spaces.
Indeed, for a generic $N \times N$ matrix of regression weights $\mbW$, the resulting distance score fails to even be symmetric:
\begin{equation}
\label{eq:naive-regression}
\min_{\mbW_i} \Vert \mbW_i \mbX_i - \mbX_j \Vert_F \neq \min_{\mbW_j} \Vert \mbX_i - \mbW_j \mbX_j \Vert_F
\end{equation}
This method of measuring similarity would therefore be sensitive to the order in which systems, such as animals or brain regions, are compared \parencite{muzellec2026reverse}.

Thus, simple forms of linear regression do not easily enable comparisons of neural systems within a metric space, as we show below.
Of course, regressions and decoding analyses remain useful, complementary tools: precisely because they target a specific variable, they can reveal whether that information is present and readable by a downstream population—a question shape distances alone cannot answer. Indeed, we use decoding throughout this paper to validate and interpret the structure that shape distances reveal.

% In our experience, the orthogonality constraint inherent to Procrustes distance is often useful to regularize the alignment transformation.
% Nonetheless, this constraint can be relaxed by computing
% \begin{equation}
% \label{eq:whitened-procrustes-distance}
% d_{\text{Linear}}(\mbX_i, \mbX_j) = d_{\text{Proc}}(\mbZ_i \mbX_i, \mbZ_j \mbX_j)
% \end{equation}
% as the distance, where $\mbZ_i$ and $\mbZ_j$ are $N \times N$ \textit{whitening matrices}.
% This avoids the problem highlighted in \cref{eq:naive-regression} while also relaxing the orthogonality constraint inherent to Procrustes distance---i.e. we have that $d_{\text{Linear}}(\mbX_i, \mbX_j) = 0$ if any invertible linear transformation brings the two neural manifolds into alignment.
% The resulting distance is a proper metric and has connections to CCA, as discussed in \textcite{Williams2021}.

% Another, more stringent, metric we will utilize is the \textit{one-to-one matching distance}, which uses a permutation of the $N$ neurons to align the responses
% \begin{equation}
% \label{eq:one-to-one-matching}
% d_{\text{Perm}}(\mbX_i, \mbX_j) = \min_{\mbP \in \mbPi(N)} \Vert \mbP \mbX_i - \mbX_j \Vert_F
% \end{equation}
% where $\mbP$ is an $N \times N$ permutation matrix.
% Note that the set of permutation matrices, $\mbPi(N)$, is a subset of all orthogonal matrices, $\cO(N)$.
% Therefore, the Procrustes distance is a less stringent measure of similarity in the sense that $d_{\text{Proc}}(\mbX_i, \mbX_j) \leq d_{\text{Perm}}(\mbX_i, \mbX_j)$.
% \textcite{khosla2024soft} describe how to extend this distance to case where the two populations consist of unequal numbers of neurons.

% In this section, we surveyed several metrics that satisfy the desirable properties of \cref{eq:metric-space-equivalence,eq:metric-space-symmetry,eq:metric-space-triangle}.
% These belong to an even broader class of distances that are referred to as \textit{generalized shape metrics} \parencite{Williams2021}. 
% Additional mathematical details are provided in the \textit{\nameref{sec:methods}}.




% If we set $\phi(\mbX) = \mbX$ (i.e. no pre-processing) and allow $\mbQ$ to be any linear transformation, then equation~(\ref{eq:generalized-euclidean-shape-metric}) reduces to least-squares linear regression.
% This is not a metric; indeed it is not even symmetric. Somewhat remarkably, if we simply constrain $\mbQ$ to be an orthogonal matrix---i.e. implementing a rotation and reflection in $N$-dimensional space---the resulting distance is a bona fide metric.
% Formally, when we set $\cG$ equal to $\cO$ (the set of $N \times N$ orthogonal matrices) and $\phi(\mbX)$ performs mean-centering of each tuning curve, then equation (\ref{eq:generalized-euclidean-shape-metric}) becomes the \textit{Procrustes size-and-shape distance}~\parencite{Gower2004}.
% Intuitively, this distance treats $\mbX_i$ and $\mbX_j$ as equivalent if they have the same ``size-and-shape'' but potentially different orientations in neural firing rate space (\hl{figure}).

% Equation~(\ref{eq:generalized-euclidean-shape-metric}) encompasses many other possibilities beyond the Procrustes distance.
% For instance, we can recover a metric that is closely related to CCA when we choose $\phi(\mbX)$ to be a whitening transformation and optimize over orthogonal matrices (i.e., $\cG = \cO$, \hl{supplement}).
% This notion of distance is invariant to more general linear deformations, whereas Procrustes distance is only invariant to rotations in neural firing rate space.
% Another useful metric for neuroscience applications is to optimize over permutations of the neurons instead of orthogonal transformations.
% This too, produces a valid metric (\hl{supplement}), which intuitively treats two neural representations as equivalent only when one can find a one-to-one matching of the neurons.

% It is important to emphasize that Equation~(\ref{eq:generalized-euclidean-shape-metric}) does not specify a single metric, but rather encompasses a broad family of metrics that quantify different aspects of representational geometry and provide complementary scientific insights.
% We refer to these as \textit{generalized shape metrics}, since they generalize classical notions of shape distance \parencite{kendall2009shape}.
% Furthermore, in \hl{Supplementary Text} we show that measures based on RSA, CKA, and CCA are all related to this framework, and some can even be viewed as special cases of shape metrics after some modification.
% However, none of these existing approaches have been formulated in a way to satisfy and leverage all three criteria of a metric space. 

